\documentclass[11pt]{article}
\usepackage[utf8]{inputenc}
\IfFileExists{lmodern.sty}{\usepackage[T1]{fontenc}\usepackage{lmodern}\usepackage{microtype}}{}
\usepackage[margin=1.05in]{geometry}
\usepackage{amsmath,amssymb,mathtools,amsthm}
\usepackage{enumitem}
\usepackage{xcolor}
\usepackage{hyperref}
\hypersetup{colorlinks=true,linkcolor=blue!45!black,citecolor=green!35!black,urlcolor=blue!55!black}
\numberwithin{equation}{section}
\newtheorem{theorem}{Theorem}[section]
\newtheorem{lemma}[theorem]{Lemma}
\newtheorem{proposition}[theorem]{Proposition}
\newtheorem{corollary}[theorem]{Corollary}
\theoremstyle{definition}
\newtheorem{definition}[theorem]{Definition}
\newtheorem{example}[theorem]{Example}
\theoremstyle{remark}
\newtheorem{remark}[theorem]{Remark}
\newcommand{\C}{\mathbb C}\newcommand{\Q}{\mathbb Q}\newcommand{\R}{\mathbb R}\newcommand{\Z}{\mathbb Z}
\newcommand{\Qbar}{\overline{\mathbb Q}}
\newcommand{\SL}{\operatorname{SL}}\newcommand{\GL}{\operatorname{GL}}\newcommand{\PGL}{\operatorname{PGL}}
\newcommand{\End}{\operatorname{End}}\newcommand{\Hom}{\operatorname{Hom}}\newcommand{\Span}{\operatorname{span}}
\newcommand{\Lie}{\operatorname{Lie}}\newcommand{\Gal}{\operatorname{Gal}}\newcommand{\wt}{\operatorname{wt}}
\newcommand{\Ncal}{\mathcal N}\newcommand{\Tcal}{\mathcal T}\newcommand{\Bcal}{\mathcal B}
\newcommand{\Hq}{H_q}
\newcommand{\dd}{\,d}

\title{Single-chain comparison for weighted-plane phases\\ for polynomial-flag exponential periods\\[4pt]
\large Corrected successor --- working note for Paper V}
\author{Working note, AI-drafted (Claude), for C.~D.~Long}
\date{6 October 2026; corrected successor}

\begin{document}
\maketitle

\noindent\textbf{Status.} Corrected successor to the frozen candidate.
The main hypotheses and conclusions are retained, including the continuation
alternative $(S^*)$. The proof includes explicit local/global transport and a
stable-lattice argument in place of general formal-exponent theory.
This working note is separate from canonical Paper V; its exact revision
and checks are recorded alongside it.

\section{Statement}

Throughout $k=\Qbar$, $k_\R=k\cap\R$, $F_0=k(z)$, $F=\C(z)$. Notation follows
Paper~V: for a polynomial phase $q$ and a compact real polynomial two-chain $D$,
\[
 I_D(\omega)(z)=\int_D e^{zq}\omega,\qquad d_z=d+z\,dq\wedge,\qquad
 \Hq=\Omega^2_{F_0[x,y]/F_0}/d_z\Omega^1_{F_0[x,y]/F_0}.
\]
$K_{\partial,0}$ is the differential field over $F_0$ generated by the
polynomial-amplitude integrals over the fixed faces of $D$ and the
exponentials at their endpoints; $K_\partial=\C K_{\partial,0}$;
$K_{\partial,\xi}$ is the field generated over $k$ by the values at $\xi$.
For a face $e$ with parametrization $\phi_e$ put $p_e=q\circ\phi_e$ and
$d_e=\deg p_e$.

\begin{definition}[The class $(W)$]\label{def:W}
Fix coprime integers $a,b\ge2$ and give $x,y$ the weights $\wt(x)=b$,
$\wt(y)=a$. A phase is of class $(W)$ if
\[
 q=\alpha x^a+\beta y^b+q_<,\qquad \alpha\beta\ne0,\qquad
 q_<\in k_\R[x,y]\ \text{of weighted degree}<ab,
\]
with $\alpha,\beta\in k_\R$. Put $\mu=(a-1)(b-1)$ and
$B=\{x^iy^j:0\le i\le a-2,\ 0\le j\le b-2\}$.
\end{definition}

The phase $y^2-x^3+3x$ of Paper~V is of class $(W)$ with $(a,b)=(3,2)$, $\mu=2$.

\paragraph{Hypotheses.} Let $q$ be of class $(W)$ and $D$ a polynomial two-chain
over $k_\R$ in the sense of Paper~V, Definition~3.1: a finite $k$-linear
combination of pushforwards of $[0,1]^2$ by maps in $k_\R[u,v]^2$. This
arithmetic restriction is used in (iii).
\begin{enumerate}[label=\textup{(\Alph*)},leftmargin=2.6em]
\item[(M)] \emph{Morse.} $q$ has $\mu$ distinct critical values
 $c_1,\dots,c_\mu$ (equivalently: all critical points are nondegenerate with
 distinct values).
\item[(Q)] \emph{Independent critical values.} $c_2-c_1,\dots,c_\mu-c_1$ are
 linearly independent over $\Q$.
\item[(F)] \emph{Faces.} For every face $e$ with $d_e\ge2$, the set
 $\{c_i-c_j:i\ne j\}$ is not contained in
 $\{u-v: u,v\ \text{critical values of }p_e\}$.
\item[(S)] \emph{Saddle.} Some real critical point $P_0$ of $q$ with
 $\det\operatorname{Hess}q(P_0)<0$ lies off the image of $\partial D$, the
 multiplicity $w$ of $D$ at $P_0$ is nonzero, and $c:=q(P_0)$ is neither the
 value of $q$ at a vertex of $D$ nor a critical value of any $p_e$ on its
 parameter interval.
\end{enumerate}
For $\mu=2$, (Q) is automatic.

\begin{theorem}[Single-chain comparison]\label{thm:main}
Let $q$ be of class $(W)$ satisfying \textup{(M)} and \textup{(Q)}, and let $D$
be a polynomial two-chain over $k_\R$ satisfying \textup{(F)} and either \textup{(S)} or \textup{(S$^*$)} below. Put
$M_m=I_D(m\dd x\wedge dy)$ for $m\in B$.
\begin{enumerate}[label=\textup{(\roman*)},leftmargin=2.4em]
\item The classes of $m\,dx\wedge dy$, $m\in B$, form an $F_0$-basis of $\Hq$,
 and a $k[z,z^{-1}]$-basis of the lattice of polynomial forms. Every
 $r\in k[x,y]$ has a unique remainder and a polynomial Stokes certificate
 \[
  r=\sum_{m\in B}a_{r,m}(z)\,m+L_xP_r+L_yQ_r,\qquad a_{r,m}\in k[z^{-1}],\quad
  P_r,Q_r\in z^{-1}k[z^{-1}][x,y],
 \]
 with $L_x=\partial_x+zq_x$, $L_y=\partial_y+zq_y$.
\item The $\mu$ functions $M_m$ are algebraically independent over $K_\partial$.
\item For every $\xi\in k^\times$ the $\mu$ numbers $M_m(\xi)$ are algebraically
 independent over $K_{\partial,\xi}$.
\item For $r\in k[x,y]$ and $\xi\in k^\times$, the value $I_D(r\dd x\wedge dy)(\xi)$
 is algebraic over $K_{\partial,\xi}$ iff it lies in $K_{\partial,\xi}$ iff
 $r\,dx\wedge dy\in(d+\xi\,dq\wedge)\Omega^1_{k[x,y]/k}$.
\end{enumerate}
\end{theorem}

\begin{remark}[Variant covering Paper V's triangle]\label{rem:Sstar}
Hypothesis (S) is used only in Lemma~\ref{lem:saddle}. It may be replaced by:
\begin{enumerate}[leftmargin=2.6em]
\item[(S$^*$)] on some interval $e_0\subset\R$ free of critical values of $q$
 and of the $p_e$, of vertex values and of $0$, the density germ $\rho$ of
 $I_D(dx\wedge dy)$ continues analytically along a path to a punctured
 neighbourhood of some $c_i$, where its variation around $c_i$ is a nonzero
 constant multiple of the vanishing period $\int_{\nu_i(t)}dx\wedge dy/dq$.
\end{enumerate}
Paper~V, Lemma~11.1 verifies (S$^*$) for its triangle (the saddle $(1,0)$ lies
outside the triangle, so (S) itself fails there). With (S$^*$), Theorem
\ref{thm:main} contains Paper~V, Theorems~1.1, 12.2 and~13.2: for that datum (F)
reads $4\ne\pm4\beta$.
\end{remark}

\begin{corollary}[Polygons]\label{cor:polygon}
Let $\min(a,b)\ge3$, let $q$ be of class $(W)$ with \textup{(M)}, \textup{(Q)},
and let $D$ be a real polygonal chain (all faces straight segments) with vertices
in $k_\R^2$. Then \textup{(F)} holds automatically. If moreover the coefficients
of $q$ and the vertices lie in a field $k_0\subset k_\R$ over which
$[k_0(c_i):k_0]\ge\max(a,b)$ for all $i$, then \textup{(S)} holds as soon as some
real saddle of $q$ lies in the interior of $D$ with nonzero multiplicity.
\end{corollary}
\begin{proof}
A straight face has $d_e\le\max(a,b)$, so $p_e$ has at most $\max(a,b)-1<\mu$
critical values and at most $(\mu-1)(\mu-2)<\mu(\mu-1)$ nonzero differences,
whereas under (Q) the $\mu(\mu-1)$ numbers $c_i-c_j$ are pairwise distinct.
Vertex values lie in $k_0$ and critical values of $p_e\in k_0[s]$ have degree at
most $d_e-1<\max(a,b)$ over $k_0$, so neither can equal $c$.
\end{proof}

\begin{example}[A rank-six instance]\label{ex:intro}
For
\[
 q=x^4+y^3+xy-y,\qquad D=\operatorname{conv}\{(0,-1),(1,-1),(\tfrac12,0)\},
\]
the six functions $\int_D x^iy^je^{zq}\dd x\dd y$ $(0\le i\le2,\ 0\le j\le1)$ are
algebraically independent over the field of the three edges, and so are their
values at every nonzero algebraic $\xi$ over the edge-value field. See
\S\ref{sec:example}.
\end{example}

\paragraph{Differential-module convention.}
For a system $Y'=EY$ we use its coefficient module: a column of basis
elements $e$ satisfies $\partial e=Ee$. Solutions are differential
homomorphisms out of this module, so the solution functor is contravariant.
Its horizontal representation is dual to the solution representation.
The natural trace duality identifies the corresponding endomorphism and
trace-free endomorphism representations. All module embeddings below are
under the covariant horizontal-fibre equivalence with these identifications.

\section{Normal form for the class \texorpdfstring{$(W)$}{(W)}}

\begin{lemma}\label{lem:normal}
Let $q$ be of class $(W)$ (coprimality is not used here).
\begin{enumerate}[label=\textup{(\alph*)},leftmargin=2.2em]
\item Every $r\in k[x,y]$ can be written
 $r=\sum_{m\in B}a_m(z)m+L_xP+L_yQ$ with $a_m\in k[z^{-1}]$ and
 $P,Q\in z^{-1}k[z^{-1}][x,y]$.
\item If $\sum_{m\in B}a_m m=L_xP+L_yQ$ with $a_m\in F_0$ and
 $P,Q\in F_0[x,y]$, then all $a_m=0$.
\item Consequently $B$ is an $F_0$-basis of $\Hq$ and a $k[z,z^{-1}]$-basis of
 $G=\Omega^2[z,z^{-1}]/d_z\Omega^1[z,z^{-1}]$. In this basis
 $\nabla_z[m]=[qm]=\sum_{m'}A_{mm'}[m']$ with
 \[
  A=A_0+A_1z^{-1}+\dots+A_sz^{-s}\in M_\mu(k[z^{-1}]),
 \]
 and $A_0$ is the matrix of multiplication by $q$ on
 $k[x,y]/(q_x,q_y)$ in the basis $B$. Under \textup{(M)} its eigenvalues are
 $c_1,\dots,c_\mu$, each simple.
\end{enumerate}
The same statements hold with $k$ replaced by $\C$, and the reduction and independence arguments also apply after specialization at every $z=\xi\ne0$.
\end{lemma}
\begin{proof}
(a) Let $x^my^n$ be a monomial of $r$, with coefficient $\gamma$, of maximal
weight among those not in $B$. If $m\ge a-1$ put $v=\gamma x^{m-a+1}y^n/(a\alpha z)$.
Then $L_xv=\gamma x^my^n+(\gamma/a\alpha)\,q_{<,x}x^{m-a+1}y^n+v_x$, and the last
two terms have weight strictly below $\wt(x^my^n)$ because
$\wt(q_{<,x})<ab-b$. Replace $r$ by $r-L_xv$. If $m<a-1$ then $n\ge b-1$ and one
uses $L_y$ symmetrically. The maximal weight of a monomial outside $B$, with its
multiplicity among the finitely many monomials of that weight, strictly
decreases, so the process stops.

(b) Among all pairs $(P,Q)$ representing the given combination choose one
minimizing $\varpi=\max(\wt P+ab-b,\ \wt Q+ab-a)$. Let $P^{\rm top}$ and $Q^{\rm top}$ be the components of $P$ and $Q$ of weights
$\varpi-ab+b$ and $\varpi-ab+a$ (one of them may be zero). The weight-$\varpi$
component of $L_xP+L_yQ$ is $z(a\alpha x^{a-1}P^{\rm top}+b\beta y^{b-1}Q^{\rm top})$. If it
is nonzero it is a nonzero element of the monomial ideal $(x^{a-1},y^{b-1})$,
which cannot equal the weight-$\varpi$ component of a combination of monomials
of $B$. Hence it vanishes, and since $x^{a-1},y^{b-1}$ is a regular sequence,
$P^{\rm top}=b\beta y^{b-1}R$, $Q^{\rm top}=-a\alpha x^{a-1}R$. As $L_x$ and $L_y$
commute, $(P-L_y(R/z),\,Q+L_x(R/z))$ represents the same element and has smaller
$\varpi$ unless $P=Q=0$. So $P=Q=0$ and all $a_m=0$.

(c) The first two assertions follow from (a) and (b). In (a), the terms of
order $z^0$ give $r\equiv\sum a_m(\infty)m$ modulo $(q_x,q_y)$, and the argument
of (b) with $z$ removed shows $B$ is a basis of the Jacobian algebra; this
identifies $A_0$. That algebra is the product of its local algebras at the
critical points, on each of which $q$ has the single eigenvalue $q(P)$, with
multiplicity the local Milnor number.
\end{proof}

\section{Period calculus and irreducibility of \texorpdfstring{$\Hq$}{Hq}}

Let $S=\C\setminus\{c_1,\dots,c_\mu\}$, $C_t=q^{-1}(t)$. The arguments below
use three geometric facts.
\begin{enumerate}[label=\textup{(G\arabic*)},leftmargin=2.8em]
\item $q$ is a locally trivial fibration over $S$; every fibre of $q$ is
 reduced and irreducible; for $t\in S$, $H_1(C_t;\C)$ has dimension $\mu$ and
 its intersection form is nondegenerate.
\item The local monodromy $T_i$ around $c_i$ is the Picard--Lefschetz
 transvection $\delta\mapsto\delta\pm\langle\delta,\nu_i\rangle\nu_i$, and the vanishing
 period $\int_{\nu_i(t)}dx\wedge dy/dq$ is holomorphic at $c_i$ with a nonzero
 value there; in particular $\nu_i\ne0$ and $T_i\ne1$.
\item $\pi_1(S)$ acts irreducibly on $H_1(C_t;\C)$.
\end{enumerate}

\begin{proposition}[Geometry of the class $(W)$]\label{prop:geometry}
For $q$ of class $(W)$ satisfying \textup{(M)}, \textup{(G1)--(G3)} hold.
\end{proposition}
\begin{proof}
\emph{Fibres.} The top weighted form of $q-t$ is $\alpha x^a+\beta y^b$, which is
irreducible in $\C[x,y]$ because $\gcd(a,b)=1$. The top form of a product is
the product of the top forms, so $q-t$ is irreducible; in particular every
fibre is reduced and irreducible.

\emph{Tameness.} Put $R=\max(|x|^{1/b},|y|^{1/a})$, $X=x/R^b$, $Y=y/R^a$. Then
$R^{-(ab-b)}q_x=a\alpha X^{a-1}+o(1)$ and $R^{-(ab-a)}q_y=b\beta Y^{b-1}+o(1)$
uniformly on $\max(|X|,|Y|)=1$ as $R\to\infty$, and one of the two leading terms
has modulus bounded below there. Hence $\|\nabla q\|\to\infty$ at infinity, so
$q$ is tame in the sense of Broughton \cite[Definition~3.1]{Br}: $\|\partial q\|$
is bounded away from $0$ outside a compact neighbourhood of the critical
points. For tame polynomials, \cite[p.~230]{Br} shows that $q$ is locally
trivial over any region of the target containing no critical value, and
\cite[Theorem~1.2 and p.~231]{Br} gives $H_1(C_t;\Z)\cong\Z^{\mu'}$ for a regular
value $t$, where $\mu'$ is the sum of the Milnor numbers at the critical points;
by Lemma~\ref{lem:normal}(c), $\mu'=\dim\C[x,y]/(q_x,q_y)=\mu$.

\emph{Local model and stability.} Write
\[
 f_\epsilon(X,Y)=\alpha X^a+\beta Y^b+
 \sum_{bi+aj<ab}q_{ij}\epsilon^{ab-bi-aj}X^iY^j.
\]
Every exponent of $\epsilon$ is a positive integer, so this is a holomorphic
deformation of $f_0=\alpha X^a+\beta Y^b$. For $\epsilon\ne0$ it equals
$\epsilon^{ab}q(\epsilon^{-b}X,\epsilon^{-a}Y)$; its critical points and values
are $(\epsilon^b x_i,\epsilon^a y_i)$ and $\epsilon^{ab}c_i$.
Choose a Milnor ball $\overline B_r$ for $f_0$, and choose $\eta>0$ so that its
levels $|t|\le2\eta$ are transverse to $\partial B_r$. Compactness and
$C^1$ convergence preserve this transversality for $f_\epsilon$ with
$|t|\le\eta$ when $|\epsilon|$ is sufficiently small. Shrink the parameter disk
so that all critical points lie in $B_{r/2}$ and all critical values have
modulus less than $\eta/2$.

At the fixed value $t_* =\eta$, the family
\[
 \{(\epsilon,X,Y):f_\epsilon(X,Y)=t_*,\ (X,Y)\in\overline B_r\}
 \longrightarrow\{\epsilon:|\epsilon|<\epsilon_0\}
\]
is a proper submersion, including on its boundary. The boundary form of
Ehresmann's theorem therefore identifies its fibres with the Milnor fibre of
$f_0$. This is also the perturbation comparison in \cite[\S2.1,
pp.~29--31]{AGV}. For a fixed small nonzero $\epsilon$, the local proper
submersion over the disk minus the critical values transports this
identification to every regular level there. Put
$S_t=f_\epsilon^{-1}(t)\cap\overline B_r$ and $C_t=f_\epsilon^{-1}(t)$.
By \cite[Theorem~2.1, p.~31]{AGV} a distinguished set of vanishing cycles
$\Delta_1,\ldots,\Delta_\mu$ is a basis of $H_1(S_t;\Z)$.
For a plane-curve singularity the kernel of the intersection form has
dimension $r_0-1$, where $r_0$ is the number of branches
\cite[\S3.6, p.~100]{AGV}. The binomial germ $f_0$ is irreducible, so this
form is nondegenerate; compare also \cite[p.~411, Example~2]{AGV}.

\emph{Compatible local and global transport.} Tameness is preserved by the
fixed nonzero rescaling. Over a relatively compact target neighborhood $U$
whose closure avoids the critical values, $\|\nabla f_\epsilon\|$ is bounded
below: use tameness outside a compact set and compactness inside it.
A bounded horizontal lift of a constant target tangent vector $v\in\C$ is
\[
 V_v=v\,\frac{\overline{\nabla f_\epsilon}}{\|\nabla f_\epsilon\|^2},
 \qquad df_\epsilon(V_v)=v.
\]
Boundary transversality supplies a smooth horizontal lift tangent to
$\partial B_r$. On a compact collar patch it to $V_v$ by a partition of unity.
The patched lift remains horizontal and bounded and is tangent to the sphere.
Its flows exist for the required finite times, preserve the sphere and its two
sides, and yield compatible local trivializations of the pair $(C_t,S_t)$.
This uses both tameness and transversality, following the bounded-lift
construction in \cite[pp.~229--230]{Br}.

Inclusion $j:S_t\hookrightarrow C_t$ preserves intersection numbers. Hence
$j_*v=0$ implies $\langle v,u\rangle=\langle j_*v,j_*u\rangle=0$ for every $u$,
and local nondegeneracy gives $v=0$. The global and local homology groups both
have dimension $\mu$, so $j_*$ is an isomorphism. The compatible transport
above makes it an isomorphism of local systems. The punctured disk contains
all critical values and its inclusion in their complement in $\C$ induces an
isomorphism on fundamental groups. Thus the global intersection form is
nondegenerate and its monodromy is that of this morsification. Rescaling back
to $q$ proves the remaining assertions of (G1).

\emph{(G2).} The Picard--Lefschetz formula is \cite[\S1.3, p.~26]{AGV}. At a
nondegenerate critical point choose holomorphic coordinates with
$q=c_i+uv$, $dx\wedge dy=J\,du\wedge dv$. The vanishing cycle is
$\{uv=t-c_i,\ |u|=r\}$ and $dx\wedge dy/dq=-J\,du/u$ on $C_t$, so the vanishing
period is $\mp2\pi i\sum_\kappa J_{\kappa\kappa}(t-c_i)^\kappa$, holomorphic with
value $\mp2\pi iJ(0)\ne0$ at $c_i$. Hence $\nu_i\ne0$, and $T_i\ne1$ by
nondegeneracy of the form.

\emph{(G3).} By the previous steps $H_1(C_t)$ has the basis
$\Delta_1,\dots,\Delta_\mu$, the global monodromy group is generated by the
Picard--Lefschetz transformations $T_{\Delta_i}$, and the form is nondegenerate.
By the theorem of Gabrielov and Lazzeri \cite[Theorem~3.3]{AGV}, the diagram of
a singularity relative to any distinguished basis is connected: the graph on
$\{1,\dots,\mu\}$ with an edge when $\langle\Delta_i,\Delta_j\rangle\ne0$ is
connected. (The diagram is defined after stabilization, which changes
off-diagonal intersection numbers only by signs; self-intersections may change.) Let $W\ne0$ be invariant. By
nondegeneracy there are $w\in W$ and $i$ with $\langle w,\Delta_i\rangle\ne0$;
then $(T_{\Delta_i}-1)w=\pm\langle w,\Delta_i\rangle\Delta_i$ puts $\Delta_i$ in
$W$, and $(T_{\Delta_j}-1)\Delta_i=\pm\langle\Delta_i,\Delta_j\rangle\Delta_j$
propagates membership along the connected graph. So $W=H_1$.
\end{proof}

For a horizontal family of cycles $\delta(t)\subset C_t$ over a simply connected
$U\subset S$ and $\omega=\sum_m z^m\omega_m\in\Omega^2_{\C[x,y]}[z]$ put
\[
 \pi_{\omega_m}(t)=\int_{\delta(t)}\frac{\omega_m}{dq},\qquad
 \Pi_\delta(\omega)(t)=\sum_m(-\partial_t)^m\pi_{\omega_m}(t).
\]
\begin{lemma}[Weyl intertwining]\label{lem:Pi}
$\Pi_\delta(z\omega)=-\partial_t\Pi_\delta(\omega)$,
$\Pi_\delta(\partial_z\omega+q\omega)=t\,\Pi_\delta(\omega)$, and
$\Pi_\delta(d_z\beta)=0$ for $\beta\in\Omega^1[z]$.
\end{lemma}
\begin{proof}
The first is the definition. Since $q=t$ on $C_t$, $\pi_{q\omega_m}=t\pi_{\omega_m}$,
and $(-\partial_t)^mt=t(-\partial_t)^m-m(-\partial_t)^{m-1}$ gives the second.
For a polynomial one-form $\beta_0$, $\pi_{dq\wedge\beta_0}=\int_{\delta(t)}\beta_0$
and $\pi_{d\beta_0}=\partial_t\int_{\delta(t)}\beta_0$ (Gelfand--Leray); the two
contributions to $\Pi_\delta((d+z\,dq\wedge)\beta)$ cancel termwise.
\end{proof}

\begin{lemma}[Period injectivity]\label{lem:injective}
Assume \textup{(G1)}, \textup{(G3)} and $\mu\ge2$. If $\Pi_\delta(\omega)\equiv0$
for every $\delta$, then $[\omega]=0$ in $\Hq\otimes\C$.
\end{lemma}
\begin{proof}
If $d\theta=\omega'$ then $d_z\theta=\omega'+z\,dq\wedge\theta$, so
$[\omega']=-z[dq\wedge\theta]$. Write $\omega=\sum_{m=0}^Nz^m\omega_m$ and apply this to the lowest coefficient:
subtracting $d_z\theta_0$ removes the $z^0$ term and changes only the
coefficient of $z^1$. After $N$ steps $\omega\equiv z^N\omega^\sharp$ modulo
$d_z\Omega^1[z]$ with $\omega^\sharp$ independent of $z$, and
$\Pi_\delta(\omega)=(-\partial_t)^N\pi_{\omega^\sharp}$ by Lemma~\ref{lem:Pi}. Thus
every branch of every period of $\alpha^\sharp=\omega^\sharp/dq$ is a polynomial of
degree $<N$ (zero when $N=0$), hence single valued: $\int_{(g-1)\delta}\alpha^\sharp=0$ for all
$g\in\pi_1(S)$. The span of the $(g-1)\delta$ is an invariant subspace, nonzero
because an irreducible representation of dimension $\ge2$ is nontrivial; by
(G3) it is all of $H_1$. So all periods of $\alpha^\sharp$ vanish. Choose
$\theta^\sharp$ with $d\theta^\sharp=\omega^\sharp$. Its periods have derivative
zero, hence are constant, hence single valued, hence (same argument) zero. So
$\theta^\sharp$ is exact along the generic fibres. The critical locus is finite
and all fibres are connected and nonempty; these give Bonnet's primitive
and quasi-fibered hypotheses for a map to the line. Thus
\cite[Definitions~1.3--1.4 and Theorem~1.5]{Bonnet} gives
$\theta^\sharp=dR+a\,dq$. Then
$\omega^\sharp=da\wedge dq=d_z(a\,dq)$ and $[\omega]=\pm z^N[\omega^\sharp]=0$.
\end{proof}

\begin{proposition}[Irreducibility]\label{prop:irred}
Assume \textup{(G1)--(G3)} and $\mu\ge2$. Every nonzero element of
$\Hq\otimes_{F_0}F$ has minimal differential operator of order $\ge\mu$ over
$F$. With Lemma~\ref{lem:normal}, $\Hq\otimes F$ is an irreducible differential
module of rank $\mu$.
\end{proposition}
\begin{proof}
Let $\eta\ne0$. Multiplying $\eta$ by an element of $F^\times$ changes neither
the differential submodule it generates nor the minimal order of an
annihilator (Leibniz), so we may assume $\eta=[\omega]$ with
$\omega\in\Omega^2_\C[z]$. Let
$P=\sum_{l=0}^rp_l(z)\nabla_z^l$, $p_l\in\C[z]$ not all zero, satisfy
$P\eta=0$. Multiplying $P$ on the left by a polynomial in $z$, which does not change $r$,
we may assume $P(z,\nabla_z)\omega=d_z\beta$ with $\beta\in\Omega^1[z]$. By
Lemma~\ref{lem:Pi}, for every $\delta$,
\[
 \widehat P\,\Pi_\delta(\omega)=0,\qquad
 \widehat P=\sum_lp_l(-\partial_t)\,t^l=\sum_{n=0}^Na_n(t)\partial_t^n,
 \qquad a_N(t)=(-1)^N\sum_l p_{l,N}\,t^l,
\]
where $N=\max\deg p_l$ and $p_{l,N}$ is the coefficient of $z^N$ in $p_l$. So
$a_N\ne0$ and $\deg a_N\le r$.
The map $\delta\mapsto\Pi_\delta(\omega)$ from $H_1(C_{t_0};\C)$ to germs at $t_0$
is linear, compatible with analytic continuation, and nonzero by
Lemma~\ref{lem:injective}; by (G3) it is injective. If $a_N(c_i)\ne0$, every
solution germ of $\widehat P$ near $c_i$ is holomorphic at $c_i$, so
continuation around $c_i$ is trivial on the image, contradicting $T_i\ne1$.
Hence $a_N$ vanishes at $c_1,\dots,c_\mu$ and $r\ge\deg a_N\ge\mu$. A proper
nonzero submodule would contain an element with minimal operator of order
$<\mu$.
\end{proof}

\section{The exponential torus and \texorpdfstring{$\SL_\mu$}{SL}}

Consider $Y'=AY$ with $A$ as in Lemma~\ref{lem:normal} and let $G\subset\GL_\mu$
be its differential Galois group over $F$, acting on the solution space $V$.

\begin{lemma}[Formal diagonalization]\label{lem:formal}
Let $A=\sum_{j\ge0}A_jw^j$, $w=1/z$, with $A_0$ having distinct eigenvalues
$c_1,\dots,c_n$. Then $Y'=AY$ has a formal fundamental matrix
$\widehat Y=\widehat H(w)\,z^{\Lambda}e^{Cz}$ with $C=\operatorname{diag}(c_i)$,
$\Lambda$ a constant diagonal matrix and $\widehat H\in\GL_n(\C[[w]])$.
\end{lemma}
\begin{proof}
After a constant change of basis $A_0=C$. Seek $\widehat H=\sum_{k\ge0}H_kw^k$,
$H_0=I$, with $\widehat H'=A\widehat H-\widehat H(C+\Lambda w)$; then
$\widehat Hz^\Lambda e^{Cz}$ is a solution. Since $(w^k)'=-kw^{k+1}$, the
coefficient of $w^k$ gives, for $k\ge1$,
\[
 -(k-1)H_{k-1}=[C,H_k]+\sum_{j=1}^kA_jH_{k-j}-H_{k-1}\Lambda .
\]
For $k=1$ put $\Lambda=\operatorname{diag}(A_1)$; the off-diagonal entries
$(c_i-c_l)(H_1)_{il}$ of $[C,H_1]$ determine the off-diagonal part of $H_1$.
For $k\ge2$ the diagonal part of the equation reads
$-(k-1)\operatorname{diag}H_{k-1}=\operatorname{diag}\bigl(A_1^{\rm off}H_{k-1}^{\rm off}
+\sum_{j=2}^kA_jH_{k-j}\bigr)$, which determines $\operatorname{diag}H_{k-1}$ from
quantities already known, and the off-diagonal part determines
$H_k^{\rm off}$. Induction on $k$ constructs $\widehat H$.
\end{proof}

\begin{proposition}\label{prop:SL}
Assume \textup{(M)}, \textup{(Q)}, \textup{(G1)--(G3)}, $\mu\ge2$. Then
$G\supseteq\SL_\mu$.
\end{proposition}
\begin{proof}
\emph{Formal solution and a diagonal subgroup.} Let $\mathfrak K=\C((w))$ with
$'=d/dz=-w^2d/dw$. By Lemma~\ref{lem:normal}(c) and Lemma~\ref{lem:formal},
$\widehat H'=A\widehat H-\widehat H(C+\Lambda w)$ with
$C=\operatorname{diag}(c_i)$, $\Lambda=\operatorname{diag}(\lambda_i)$,
$\widehat H\in\GL_\mu(\C[[w]])$. Let $K_0$ be a Picard--Vessiot extension of
$\mathfrak K$ for the diagonal system $Y'=(C+\Lambda w)Y$; it is generated by
nonzero $u_i$ with $u_i'=(c_i+\lambda_iw)u_i$. Then
$\widehat Y=\widehat H\operatorname{diag}(u_i)$ is a fundamental matrix of
$Y'=AY$ in $K_0$, and since $K_0$ has constants $\C$,
$\widehat L=F(\widehat Y_{ij})$ is a Picard--Vessiot extension of $F$; $G$ is its
group, written in the basis of columns of $\widehat Y$. Every
$\sigma\in\Gamma=\Gal(K_0/\mathfrak K)$ multiplies each $u_i$ by a constant, so
it restricts to an element of $G$ which is diagonal. Let $D_\Gamma\subset G$ be
the resulting closed subgroup of the diagonal torus. A character
$x^n=\prod x_i^{n_i}$ is trivial on $D_\Gamma$ iff $\prod u_i^{n_i}$ is
$\Gamma$-fixed, i.e.\ lies in $\mathfrak K$. An element $f\in\mathfrak K^\times$
has $f'/f\in w\C[[w]]$, while
$(\prod u_i^{n_i})'/\prod u_i^{n_i}=\sum n_ic_i+(\sum n_i\lambda_i)w$; so such
$n$ satisfy $\sum n_ic_i=0$. A closed subgroup of a torus is the common kernel
of the characters trivial on it, hence
\[
 D_\Gamma\supseteq T:=\{d:\ d^n=1\ \text{whenever}\ \textstyle\sum n_ic_i=0\}.
\]
$T$ is a torus with character group $\Z^\mu/\{n:\sum n_ic_i=0\}\cong\sum\Z c_i$,
acting on the $i$th basis vector with weight $c_i$; being connected,
$T\subset G^0$.

\emph{Root decomposition.} By (Q) the weights $c_i-c_j$ $(i\ne j)$ of $T$ on
$\mathfrak{gl}_\mu$ are pairwise distinct and nonzero, so
$\mathfrak g=\Lie G^0=(\mathfrak g\cap\mathfrak d)\oplus\bigoplus_{(i,j)\in R}\C E_{ij}$,
$\mathfrak d$ the diagonal matrices. $G$ is reductive because $V$ is a faithful
irreducible representation (Proposition~\ref{prop:irred}), so $G^0$ is
reductive; $\mathfrak g\cap\mathfrak d$ is the centralizer of $\Lie T$, a Cartan
subalgebra, and therefore $R=-R$. Also $R$ is closed under
$(i,j),(j,l)\mapsto(i,l)$. So $R$ defines an equivalence relation with classes
$I_1,\dots,I_s$; each $V_{I_a}=\Span\{e_i:i\in I_a\}$ is an irreducible
$G^0$-module and they are pairwise nonisomorphic, so the $G^0$-submodules of
$V$ are exactly the sums of the $V_{I_a}$.

\emph{One class.} Since $G^0$ is normal, $G$ permutes the $V_{I_a}$. The fixed
field $E_0$ of $G^0$ in $\widehat L$ is a finite Galois extension of $F$ with
group $G/G^0$. Under the isomorphism of $\widehat L$ with the Picard--Vessiot
field generated by an analytic fundamental matrix, its elements are algebraic
functions meromorphic on the universal cover of $\C^\times$ (the system is
singular only at $0,\infty$), so $E_0=F(\zeta)$ with $\zeta^d=z$. The image of
$D_\Gamma$ in $\Gal(E_0/F)$ has fixed field $E_0\cap\mathfrak K$. This is
$F(\zeta^{d/d'})$ for some $d'\mid d$, and
$(\zeta^{d/d'})^{d'}=w^{-1}$ forces $d'=1$ because $w$-orders in $\mathfrak K$
are integers. So $D_\Gamma$ maps onto $G/G^0$, i.e.\ $G=G^0D_\Gamma$. Diagonal
matrices preserve each $V_{I_a}$, so each $V_{I_a}$ is $G$-stable and
irreducibility gives $s=1$. Then every $E_{ij}$ lies in $\mathfrak g$ and
$\SL_\mu\subset G^0$.
\end{proof}

\section{Face disjointness}

\paragraph{Stable lattices at infinity.}
Put $R_\infty=\C[[w]]$, $K_\infty=\C((w))$ and
$\delta=d/dz=-w^2d/dw$. A lattice in a differential module over $K_\infty$
is a finite free $R_\infty$-submodule spanning it. A lattice $L$ is stable if
$\partial L\subset L$. Since $\delta(R_\infty)\subset wR_\infty$, its
operator induces a $\C$-linear operator on $L/wL$, called the leading
operator of this stable lattice. In the coefficient-module convention its
matrix for a system $Y'=E(w)Y$ is $E(0)^{\mathsf T}$, which has the same
spectrum as $E(0)$.

\begin{lemma}[Stable lattices and subquotients]\label{lem:lattice}
Let $L\subset M$ be a stable lattice and $N\subset M$ a differential
submodule. Then $L_N=L\cap N$ and $L_{M/N}=\operatorname{im}(L\to M/N)$
are stable lattices, and
\[
 0\longrightarrow L_N/wL_N\longrightarrow L/wL
 \longrightarrow L_{M/N}/wL_{M/N}\longrightarrow0
\]
is exact and compatible with the leading operators. Their characteristic
polynomials therefore multiply. In particular a direct sum of successive
quotients of a differential-module filtration has a stable lattice with the
same leading characteristic polynomial as $L$.
\end{lemma}
\begin{proof}
The intersection is a saturated, finitely generated submodule of $L$:
if $wv\in L\cap N$ and $v\in L$, then $v\in N$ because $N$ is a
$K_\infty$-vector space. The intersection and the quotient are free over the
discrete valuation ring $R_\infty$ and span the desired spaces. Stability
is inherited. Since the quotient is free, reduction modulo $w$ preserves
exactness. Also $\delta(w)=-w^2$, so each leading operator is well-defined
and the sequence intertwines them. Apply this successively to the filtration.
\end{proof}

\begin{lemma}[Detection of a rank-one leading coefficient]\label{lem:rankone}
If a module $M$ with stable lattice $L$ contains a nonzero vector $v$ with
$\partial v=(d+\rho w)v$, $d,\rho\in\C$, then $d$ is an eigenvalue of the
leading operator on $L/wL$.
\end{lemma}
\begin{proof}
The saturated lattice $L\cap K_\infty v$ has a generator $fv$, where
$f=w^m u$ with $u\in R_\infty^\times$. Its coefficient is
\[
 d+\rho w+\frac{\delta f}{f},\qquad
 \frac{\delta f}{f}=-mw-w^2\frac{u_w}{u}\in wR_\infty.
\]
Its leading eigenvalue is therefore $d$. Its reduction injects into $L/wL$
by Lemma~\ref{lem:lattice}.
\end{proof}

\begin{corollary}[A leading-spectrum obstruction]\label{cor:face-lattice}
Suppose $M$ has a stable lattice with leading spectrum $\Sigma$. Let $P$
admit a basis over $K_\infty$ with
$\partial e_i=(c_i+\lambda_iw)e_i$. If
\[
 \End^0(P)\hookrightarrow\End(M^{\rm ss})
\]
is a differential-module embedding, then $c_i-c_j\in\Sigma-\Sigma$ for
all $i\ne j$.
\end{corollary}
\begin{proof}
By Lemma~\ref{lem:lattice}, $M^{\rm ss}$ has a stable lattice $L^{\rm ss}$
with the same leading characteristic polynomial. The lattice
$\End_{R_\infty}(L^{\rm ss})$ is stable; its leading operator is the
commutator with that of $L^{\rm ss}$, with eigenvalues $s-s'$ for
$s,s'\in\Sigma$. The vector $e_i\otimes e_j^*$ is trace-free for $i\ne j$
and has differential coefficient
$c_i-c_j+(\lambda_i-\lambda_j)w$. Its image in the endomorphism module is
nonzero, so Lemma~\ref{lem:rankone} gives the result.

If $M^{\rm ss}$ was taken over $F$ first, extend a composition series from
$F$ to $K_\infty$ and apply the lattice lemma to this filtration. Its
quotients need not remain simple. No assertion that semisimplification
commutes with this extension is used.
\end{proof}

For a face polynomial $p$ of degree $d\ge2$, the one-variable reduction gives
a matrix in $M_{d-1}(k[z^{-1}])$ whose leading matrix is multiplication by
$p$ on $k[s]/(p')$. Thus its stable lattice has as leading spectrum the
critical values of $p$, also when those values repeat. For the interior
module, Lemma~\ref{lem:formal} supplies the diagonal basis required in
Corollary~\ref{cor:face-lattice}. These are all the local facts needed below;
no general formal-normal-form theorem is required.

Let $\Bcal$ be the closed homogeneous system satisfied by the boundary vector
(reduced face moments, endpoint exponentials, $1$), $\Tcal_\partial$ its tensor
category and $G_\partial$ its Galois group.

\begin{proposition}\label{prop:faces}
Suppose $G\supseteq\SL_\mu$, $\mu\ge2$, and that the interior module
has the formal diagonalization of Lemma~\ref{lem:formal}. If $\End^0(A)$
belongs to $\Tcal_\partial$, then there is a face $e$ with $d_e\ge2$ such that
every $c_i-c_j$ is a difference of two critical values of $p_e$. In particular
\textup{(F)} implies $\End^0(A)\notin\Tcal_\partial$.
\end{proposition}
\begin{proof}
The Galois group of $\End^0(A)$ is $\PGL_\mu$, so there is a surjection
$\rho:G_\partial\to\PGL_\mu$ through which $G_\partial$ acts on
$\mathfrak{sl}_\mu$ by $\operatorname{Ad}\circ\rho$. The unipotent radical dies,
so $\rho$ factors through $G'=\Gal(\Bcal^{\rm ss})\subset\prod_e\GL(W_e)\times(\text{torus})$,
$W_e$ the solution space of $H_{p_e}^{\rm ss}$. Write
$\Lie G'=\mathfrak z\oplus\mathfrak s_1\oplus\dots\oplus\mathfrak s_r$. Since
$\mathfrak{sl}_\mu$ is simple, $\ker d\rho$ contains $\mathfrak z$ and all simple
ideals but one, $\mathfrak s$; as $\ker d\rho$ is $\operatorname{Ad}(G')$-stable,
so is $\mathfrak s$, and $d\rho:\mathfrak s\to\mathfrak{sl}_\mu$ is a
$G'$-equivariant isomorphism. The projections
$\mathfrak s\to\mathfrak{gl}(W_e)$ are $G'$-equivariant and each is injective or
zero; they are not all zero because the remaining factor is abelian. For an
$e$ with injective projection, the $G_\partial$-module $\mathfrak{sl}_\mu$ embeds
in $\End(W_e)$. Using the covariant horizontal-fibre equivalence and the trace dualities
noted above, this gives a differential-module embedding
\[
 \End^0(A)\hookrightarrow\End(H_{p_e}^{\rm ss}).
\]
For the face module the leading spectrum is the set of critical values of
$p_e$. Corollary~\ref{cor:face-lattice}, after extension to $K_\infty$,
therefore makes every $c_i-c_j$ a difference of two such critical values.
This contradicts (F).
\end{proof}

\section{The saddle lemma}

Let $\Ncal=\Span_F\{1,\ \text{reduced face moments},\ \text{endpoint exponentials}\}$.

\begin{lemma}\label{lem:saddle}
Assume \textup{(G1)--(G3)}, $\mu\ge2$, and \textup{(S)} or \textup{(S$^*$)}. Then
$I_D(dx\wedge dy)\notin\Ncal$.
\end{lemma}
\begin{proof}
Suppose $P(z)I_D(dx\wedge dy)=\sum_jR_j(z)g_j$ with $0\ne P$, $R_j\in\C[z]$ and
$g_j$ among the generators. As in Paper~V, Lemma~11.1, the inverse Laplace
transforms of the $g_j$ have algebraic densities on intervals free of the
finitely many face-critical, vertex and zero values, and point masses elsewhere.
The compactly supported distributions are determined by their Laplace
transforms: restrict an identically zero entire transform to the imaginary
axis and apply Fourier uniqueness. Hence on such an interval the density
$\rho$ of $I_D(dx\wedge dy)$ satisfies:
$g=P(-\partial_t)\rho$ is algebraic over $\C(t)$. Under (S) this applies to
$(c,c+\varepsilon)$ for small $\varepsilon>0$.

\emph{Local form under (S).} By the real-analytic Morse lemma there are local
coordinates with $q=c+uv$, $dx\wedge dy=J\,du\wedge dv$,
$J(0)^2=|\det\operatorname{Hess}q(P_0)|^{-1}$. Take a box
$\{|u|,|v|\le\delta\}$ on which the multiplicity of $D$ is $w$. After shrinking $\delta$ the edges of the box are transverse to the levels near
$c$ and its corners have levels $c\pm\delta^2\ne c$. Outside the open
box the level curves near level $c$ meet neither a critical point (distinct
critical values) nor a tangency or vertex of $\partial D$ (by (S)), so that
contribution to $\rho$ is analytic at $c$. Inside, with $s=t-c$,
\[
 \rho_{\rm box}(s)=\pm w\int_{|s|/\delta\le|u|\le\delta}J(u,s/u)\frac{du}{|u|}
 =\pm w\bigl(-2j(s)\log s+\text{analytic}\bigr)\quad(0<s<\delta^2),
\]
where $J=\sum J_{\kappa\lambda}u^\kappa v^\lambda$ and $j(s)=\sum_\kappa J_{\kappa\kappa}s^\kappa
=\frac1{2\pi i}\oint_{|u|=r}J(u,s/u)\frac{du}u$; the terms with
$\kappa\ne\lambda$ integrate to convergent power series in $s$. Since
$dx\wedge dy/dq=-J\,du/u$ on $C_t$, $j$ is $\mp\frac1{2\pi i}$ times the period
$h(t)=\int_{\nu(t)}dx\wedge dy/dq$ over the vanishing cycle, and $j(0)=J(0)\ne0$.
Thus $\rho$ continues around $c$ with variation a nonzero constant times $h$,
which is (S$^*$).

\emph{Conclusion from (S$^*$).} Write the variation as $\lambda h$, $\lambda\ne0$. The
vanishing period $h$ is fixed by the local monodromy, so $N$ turns around
$c_i$ change $\rho$ by $N\lambda h$. The germ $g$ continues with $\rho$ and
remains algebraic, so for some $N\ge1$ it returns to itself; therefore
$N\lambda P(-\partial_t)h=0$ near $c_i$ and $h$ extends to an entire function. So $h$ is
single valued on $S$, i.e.\ $\int_{(g-1)\nu_i}dx\wedge dy/dq=0$ for all
$g\in\pi_1(S)$. As in Lemma~\ref{lem:injective} the $(g-1)\nu_i$ span $H_1$
($\nu_i\ne0$ is not invariant, by (G3) and $\mu\ge2$), so the period over $\nu_i$
itself vanishes, contradicting (G2).
\end{proof}

\section{Proof of Theorem \ref{thm:main}}

(i) is Lemma~\ref{lem:normal}. For (ii): integrating the certificates of
$qm$ gives $M'=AM+b$ with $b_m=\int_{\partial D}e^{zq}(P_{qm}\,dy-Q_{qm}\,dx)\in\Ncal$.
Paper~V, Theorem~8.2 (the actual one-column criterion) requires:
(1) $G\supseteq\SL_\mu$, which is Proposition~\ref{prop:SL};
(2) $\End^0(A)\notin\Tcal_\partial$, which is Proposition~\ref{prop:faces}
with (F) and Lemma~\ref{lem:formal};
(3) some $M_m\notin\Ncal$, which is Lemma~\ref{lem:saddle} since $1\in B$.
Hence the $M_m$ are algebraically independent over $F(\Ncal)=K_\partial$.

(iii) The vector (boundary coordinates, $M$) consists of $E$-functions
(Paper~V, Theorem~2.1) and satisfies a homogeneous linear system over $k(z)$
whose only finite pole is $z=0$: for $M$ by Lemma~\ref{lem:normal}, for the
faces by its one-variable case. Let $d$ be the transcendence degree of
$K_{\partial,0}$ over $k(z)$. Part (ii) gives transcendence degree $d+\mu$
for the combined vector. Beukers' form of Siegel--Shidlovskii
\cite[Theorem~1.1, p.~370]{Beukers}, applied first to the boundary vector
and then to the combined vector, gives the numerical degrees $d$ and
$d+\mu$ at every $\xi\in k^\times$. Their difference is $\mu$.
The one-variable normal forms have only powers of $z$ as denominators, so
these reduced boundary values generate the stated face-value field at
all such $\xi$.

(iv) Specialize (i) at $z=\xi$: the boundary term lies in $K_{\partial,\xi}$,
so algebraicity forces all $a_{r,m}(\xi)=0$ by (iii), which gives the exactness;
conversely exactness gives membership by Stokes. \qed

\section{The rank-six example}\label{sec:example}

Let $q=x^4+y^3+xy-y$, so $(a,b)=(4,3)$, $\mu=6$,
$B=\{1,x,x^2,y,xy,x^2y\}$, and $D$ the triangle of Example~\ref{ex:intro}.
The following are exact computations (script supplied).
\begin{itemize}[leftmargin=1.6em]
\item \emph{Connection.} $A=A_0+A_1/z$ with
 $\det(t-A_0)$ proportional to
 \[
 \chi(t)=1289945088\,t^6-573308928\,t^4+83980800\,t^3+83854656\,t^2+5477429\,t-4197429.
 \]
 The same polynomial is obtained independently as the resultant eliminating
 the critical points.
\item \emph{(M) and (Q).} $\chi$ is irreducible over $\Q$; modulo $7$ it
 factors with degrees $(1,5)$ and modulo $31$ with degrees $(1,2,3)$. A
 transitive group of degree $6$ containing a $5$-cycle is primitive, and with
 a transposition (the cube of the second Frobenius) it is $S_6$. Given a
 relation $\sum a_i c_i=0$ with rational $a_i$, applying a transposition
 and subtracting yields $(a_i-a_j)(c_i-c_j)=0$. Thus all $a_i$ are equal;
 a relation among the differences has coefficient sum zero and is trivial.
\item \emph{(S).} The critical points are $x=1-3y^2$ with
 $-108y^6+108y^4-36y^2+y+4=0$. Write $f(y)$ for this sextic. Sturm counts give exactly two real roots, one in
 $(-\tfrac{420}{1000},-\tfrac{419}{1000})$ and one in
 $(\tfrac{722}{1000},\tfrac{723}{1000})$. On the critical locus
 $\det\operatorname{Hess}q=72x^2y-1=-f'(y)$, and $f'$ has no root in either
 interval, with $f'>0$ on the first and $f'<0$ on the second. So the first root
 gives a saddle $P_0\approx(0.471644,-0.419665)$, $c=q(P_0)\approx0.197305$,
 and the second does not. The polynomials $x-\tfrac{y+1}2$ and
 $1-\tfrac{y+1}2-x$ with $x=1-3y^2$ have no root on the first interval and are
 positive there, so $P_0$ is interior to $D$.
\item \emph{(F) and the rest of (S)} hold by Corollary~\ref{cor:polygon} with
 $k_0=\Q$, since $[\Q(c):\Q]=6\ge4$.
\end{itemize}
For reproducibility, writing $\bar\chi_p$ for the monic reduction modulo $p$,
\begin{align*}
 \bar\chi_5&=t^6-t^4+2t^2-2t+2,\\
 \bar\chi_7&=(t+3)(t^5-3t^4+3t^2+3t+2),\\
 \bar\chi_{31}&=(t-14)(t^2+8t+10)(t^3+6t^2-14t+1).
\end{align*}
The displayed factors modulo $7,31$ are irreducible and squarefree.
Irreducibility modulo $5$ is certified by
\[
 \gcd(\bar\chi_5,t^{5^2}-t)=\gcd(\bar\chi_5,t^{5^3}-t)=1,
 \qquad t^{5^6}\equiv t\pmod{\bar\chi_5}.
\]
All these are exact polynomial computations in the accompanying certificate.
The supplementary exact checks give
$\operatorname{tr}(A_0^k(A_1+I))=0$ for $0\le k\le5$.
Since $A_0$ has simple spectrum, the Vandermonde matrix in its eigenvalues
shows that all six diagonal entries of $A_1$ in its eigenbasis are $-1$.
Thus the six powers in Lemma~\ref{lem:formal} are exactly $-1$.

Numerical sanity checks, not part of the proof: the Stokes identities for
$q\cdot m$ hold on $D$ to $10^{-17}$ by quadrature at $z=0.8$, and the density
has $\rho(c\pm s)\approx-0.7197\log s$ on both sides, matching
$2/\sqrt{7.7215}=0.71975$.

\section{External inputs and scope}\label{sec:inputs}
The proof uses the following precise external results. The elementary
formal diagonalization and stable-lattice lemmas are proved here; neither
Wasow's formal simplification theorem nor a general formal normal form is
an input.
\begin{enumerate}[leftmargin=2.2em]
\item Broughton \cite{Br}, Definition~3.1, p.~225; Theorem~1.2, p.~219,
 and the curve calculation p.~231; bounded horizontal lifts and local
 triviality, pp.~229--230. Proposition~3.4, p.~228, gives an unused
 Newton-theoretic alternative to our explicit tameness estimate.
\item Arnold--Gusein-Zade--Varchenko \cite{AGV}, \S1.3, p.~26
 (Picard--Lefschetz); \S2.1, pp.~29--31 and Theorem~2.1
 (perturbed fibre and distinguished basis); Theorem~3.3, p.~75
 (connectedness); \S3.6, p.~100 and p.~411, Example~2
 (nondegeneracy). Proposition~\ref{prop:geometry} includes the parameter
 choices and the proper-submersion argument for the perturbation comparison.
 Connectedness can alternatively be read in \cite[Theorem~2, p.~274]{La};
 \cite[Corollary~4, p.~27]{Ga} gives an explicit grid basis, and
 \cite[Theorem~17, Corollary~18]{Eb} explains the general statement.
\item Bonnet \cite[Definitions~1.3--1.4 and Theorem~1.5, p.~375]{Bonnet},
 for polynomial relative exactness; Singer \cite[Theorem~1.5.2]{Singer},
 for the Picard--Vessiot tensor equivalence; Beukers
 \cite[Theorem~1.1, p.~370]{Beukers}, for specialization of
 transcendence degrees.
\item The polynomial boundary-system theorem and actual one-column
 criterion of the pinned version of Paper~V \cite[Theorems~2.1 and~8.2]{PaperV}.
 Its Lemma~11.1 supplies the continuation verification for the elliptic
 triangle. These inputs do not use the present general theorem.
\end{enumerate}

\paragraph{Scope.} The fields here belong to the finitely many fixed faces,
not to all compact line periods. The theorem concerns one actual chain and
one phase. It does not prove a several-chain theorem, a several-phase
comparison, or general formal completeness. Hypothesis (Q) is sufficient;
its failure alone does not imply a smaller differential Galois group.
Genericity of (Q) in class $(W)$ is not established. No such genericity is
needed for the theorem under its explicit hypotheses.
For $\min(a,b)=2$ the face condition must still be checked; for example it
fails on horizontal edges when $q=y^2+f(x)$.

\paragraph{Revision record.} This is a corrected successor to the immutable
6 October candidate. It incorporates the geometric transport details,
the coefficient-module convention, a stable-lattice replacement for the
formal-exponent argument, precise references, and the off-diagonal
stabilization and (Q)-scope corrections. It does not weaken (Q), expand
the class of chains, or introduce a formal-completeness conclusion.
The computational certificates establish the displayed finite example;
they do not prove the general geometric or Galois assertions.

\begin{thebibliography}{99}
\bibitem{AGV} V.~I.~Arnold, S.~M.~Gusein-Zade, A.~N.~Varchenko,
 \emph{Singularities of Differentiable Maps}, vol.~II, Birkh\"auser, 1988.
\bibitem{Br} S.~A.~Broughton, \emph{Milnor numbers and the topology of
 polynomial hypersurfaces}, Invent.\ Math.\ \textbf{92} (1988), 217--241.
\bibitem{Eb} W.~Ebeling, \emph{Distinguished bases and monodromy of complex
 hypersurface singularities}, arXiv:1905.12435.
\bibitem{Ga} A.~M.~Gabrielov, \emph{Intersection matrices for certain
 singularities}, Funktsional.\ Anal.\ i Prilozhen.\ \textbf{7} (1973), no.~3,
 18--32.
\bibitem{La} F.~Lazzeri, \emph{A theorem on the monodromy of isolated
 singularities}, Ast\'erisque \textbf{7--8} (1973), 269--275.

\bibitem{Bonnet} P.~Bonnet,
 \emph{Relative exactness modulo a polynomial map and algebraic
 $(\mathbb C^p,+)$-actions}, Bull. Soc. Math. France \textbf{131} (2003),
 no.~3, 373--398. \href{https://doi.org/10.24033/bsmf.2447}{doi:10.24033/bsmf.2447}.
\bibitem{Beukers} F.~Beukers,
 \emph{A refined version of the Siegel--Shidlovskii theorem}, Ann. of Math.
 \textbf{163} (2006), 369--379.
 \href{https://doi.org/10.4007/annals.2006.163.369}{doi:10.4007/annals.2006.163.369}.
\bibitem{Singer} M.~F.~Singer,
 \emph{Introduction to the Galois theory of linear differential equations},
 \href{https://arxiv.org/abs/0712.4124}{arXiv:0712.4124v2}, 2008.
 We use ordinary Picard--Vessiot theory and Theorem~1.5.2, not the
 general formal-normal-form theorem of \S1.4.
\bibitem{PaperV} C.~D.~Long and A.-A.~Le~Blanc,
 \emph{Relative Exponential Periods on Polynomial Flags:
 Polynomial Boundary Systems and an Elliptic Comparison Theorem},
 Paper~V, September 2026; repository snapshot
 \href{https://github.com/long-mathematics/exponential-integral-theorem/tree/d5ae8f9081c7bc34d1622316bd8463482f244395}{d5ae8f90},
 inspected 6 October 2026.

\end{thebibliography}

\end{document}
